% Session 4 -- why Landau levels are flat: the Landau-gauge reduction turns every
% k_y into the same harmonic oscillator, only displaced.
% Placement: sec. "Landau levels of Schroedinger electrons", after equation (4.6).
%
% LAYOUT NOTE: two pgfplots axes in one tikzpicture; the second is placed off the
% first with "at={($(ax1.south east)+(...)$)}". Everything inside each axis is in
% axis cs, nothing is drawn in bare tikzpicture coordinates.
\begin{aqmfigure}{fig:s04:guidingcenter}{Why the Landau levels are flat. (a) In the Landau gauge each
$k_y$ reduces (4.5) to a one-dimensional harmonic oscillator in $x$, identical in every respect
except that it is centred on the guiding centre $x_0=\hbar k_y/B$: raising $k_y$ slides the whole
well to the left without changing its curvature. Every well therefore carries the identical ladder
$E_n=\hbar\omega_c(n+\tfrac12)$, drawn here as the horizontal guides, and the ground-state
wavepacket of each one is the same shape parked in a different place. (b) Plotted against $k_y$,
that is a set of exactly dispersionless bands. The flatness is not an artefact of the gauge: it is
the statement that the energy knows nothing about where in the sample the electron sits, which is
also where the enormous Landau-level degeneracy comes from.}

\begin{tikzpicture}

% ============================ (a) the shifted wells ==========================
\begin{axis}[
  aqmplot,
  name=ax1,
  width=8.8cm, height=5.8cm,
  axis lines=left,
  axis line style={aqm, black},
  title={(a) the Landau-gauge oscillator, one per $k_y$},
  xlabel={$x$}, ylabel={$E$},
  xlabel style={at={(axis description cs:1.02,0)}, anchor=west},
  xmin=-6.2, xmax=6.2, ymin=0, ymax=4.6,
  xtick=\empty,
  ytick={0.5,1.5,2.5}, yticklabels={$E_0$,$E_1$,$E_2$},
  enlargelimits=false,
]

% --- the shared ladder: the whole point of the figure ------------------------
\addplot[aqmguide, forget plot] coordinates {(-6.0,0.5) (6.0,0.5)};
\addplot[aqmguide, forget plot] coordinates {(-6.0,1.5) (6.0,1.5)};
\addplot[aqmguide, forget plot] coordinates {(-6.0,2.5) (6.0,2.5)};

% --- the three displaced parabolic wells -------------------------------------
% muted+dashed for the two flanking ones, accent for the central one.
\addplot[aqmcurvemuted, domain=-5.50:-0.90, samples=60, forget plot] {0.55*(x+3.2)^2};
\addplot[aqmcurvemuted, domain= 0.90: 5.50, samples=60, forget plot] {0.55*(x-3.2)^2};
\addplot[aqmaccent, line width=1.1pt, domain=-2.30:2.30, samples=60, forget plot] {0.55*x^2};

% --- their (identical) ground-state wavepackets ------------------------------
\addplot[aqmmuted, line width=0.8pt, domain=-4.9:-1.5, samples=60, forget plot]
  {0.5+0.52*exp(-2*(x+3.2)^2)};
\addplot[aqmmuted, line width=0.8pt, domain=1.5:4.9, samples=60, forget plot]
  {0.5+0.52*exp(-2*(x-3.2)^2)};
\addplot[aqmaccent, line width=1.1pt, domain=-1.7:1.7, samples=60, forget plot]
  {0.5+0.52*exp(-2*x^2)};

% --- guiding centres marked on the axis --------------------------------------
\node[circle, fill=black, inner sep=0pt, minimum size=3.4pt] at (axis cs:-3.2,0) {};
\node[circle, fill=black, inner sep=0pt, minimum size=3.4pt] at (axis cs: 0,0)   {};
\node[circle, fill=black, inner sep=0pt, minimum size=3.4pt] at (axis cs: 3.2,0) {};
\node[aqmlabsmall, anchor=north] at (axis cs:0,-0.20) {guiding centres $x_0=\hbar k_y/B$};

% --- the sliding statement ---------------------------------------------------
\draw[aqmarrowaccent] (axis cs:2.6,3.45) -- (axis cs:-2.6,3.45);
\node[aqmlabsmall, anchor=south] at (axis cs:0,3.52) {increasing $k_y$};
\node[aqmlabsmall, anchor=south] at (axis cs:0,4.02)
  {\itshape same ladder, different place};

\end{axis}

% ============================== (b) the flat band ============================
\begin{axis}[
  aqmplot,
  name=ax2,
  at={($(ax1.south east)+(2.1cm,0)$)}, anchor=south west,
  width=4.4cm, height=5.8cm,
  axis lines=left,
  axis line style={aqm, black},
  title={(b) the band},
  xlabel={$k_y$}, ylabel={$E$},
  xlabel style={at={(axis description cs:1.02,0)}, anchor=west},
  xmin=-1.15, xmax=1.15, ymin=0, ymax=4.6,
  xtick={0}, xticklabels={$0$},
  ytick={0.5,1.5,2.5}, yticklabels={$E_0$,$E_1$,$E_2$},
  enlargelimits=false,
]
\addplot[black, line width=1.1pt, forget plot] coordinates {(-1.05,0.5) (1.05,0.5)};
\addplot[black, line width=1.1pt, forget plot] coordinates {(-1.05,1.5) (1.05,1.5)};
\addplot[black, line width=1.1pt, forget plot] coordinates {(-1.05,2.5) (1.05,2.5)};
\node[circle, fill=aqmaccent, inner sep=0pt, minimum size=4pt] at (axis cs:-0.55,0.5) {};
\node[circle, fill=aqmaccent, inner sep=0pt, minimum size=4pt] at (axis cs: 0,0.5)    {};
\node[circle, fill=aqmaccent, inner sep=0pt, minimum size=4pt] at (axis cs: 0.55,0.5) {};
\node[aqmlabsmall, anchor=south, align=center] at (axis cs:0,3.05)
  {$\partial E/\partial k_y=0$\\[-1pt]\itshape exactly flat};
\end{axis}

% =========================== the degeneracy punchline ========================
\node[aqmboxaccent, anchor=north west, align=center]
  at ($(ax1.south west)+(0,-1.30cm)$)
  {on a sample of size $W\times L_y$ the allowed $k_y$ are spaced by $2\pi/L_y$,\\[1pt]
   so the guiding centres are spaced by $\Delta x_0=2\pi\hbar/BL_y$\\[4pt]
   number of orbitals per Landau level $=W/\Delta x_0=BWL_y/2\pi\hbar=\Phi/\Phi_0$:\\[1pt]
   one state per flux quantum through the sample};

\end{tikzpicture}

\end{aqmfigure}
